\documentclass[12pt]{article}

\usepackage{amsmath,amssymb,amsthm}
\usepackage{tikz-cd}
\usepackage{hyperref}
\usepackage{cleveref}
\usepackage{graphicx}
\usepackage{tikz}
\usepackage{everypage}
\usepackage{xcolor}
\usepackage[margin=1in]{geometry}
\usepackage{enumitem}

\hypersetup{
  colorlinks=true,
  linkcolor=blue!50!black,
  citecolor=blue!50!black,
  urlcolor=blue!50!black
}

\definecolor{grokgray}{RGB}{110,110,110}
\AddEverypageHook{%
  \ifnum\value{page}=1
    \begin{tikzpicture}[remember picture, overlay]
      \node[
        rotate=90,
        anchor=south,
        font=\Large\sffamily\bfseries\color{grokgray},
        inner sep=0pt
      ] at ([xshift=38pt, yshift=0.52\paperheight]current page.south west)
      {GrokRxiv:2026.05.paper-04-rational-dilation-externalization\quad
       [\,math.CT\,]\quad
       07 May 2026};
    \end{tikzpicture}
  \fi
}

\theoremstyle{plain}
\newtheorem{theorem}{Theorem}[section]
\newtheorem{proposition}[theorem]{Proposition}
\newtheorem{lemma}[theorem]{Lemma}
\newtheorem{corollary}[theorem]{Corollary}

\theoremstyle{definition}
\newtheorem{definition}[theorem]{Definition}
\newtheorem{example}[theorem]{Example}

\theoremstyle{remark}
\newtheorem{remark}[theorem]{Remark}

% Notation
\newcommand{\HH}{\mathcal{H}}
\newcommand{\KB}{K_{B}}
\newcommand{\HardyTwo}{H^{2}}
\newcommand{\Disk}{\mathbb{D}}
\newcommand{\Real}{\mathbb{R}}
\newcommand{\Complex}{\mathbb{C}}
\newcommand{\Rationals}{\mathbb{Q}}
\newcommand{\Naturals}{\mathbb{N}}
\newcommand{\Integers}{\mathbb{Z}}
\newcommand{\fracpart}[1]{\{#1\}}
\newcommand{\Yon}{\mathrm{y}}
\newcommand{\Mellin}{\mathcal{M}}
\newcommand{\PSh}{\mathbf{PSh}}
\newcommand{\Set}{\mathbf{Set}}
\newcommand{\DQ}{D_{Q}}
\newcommand{\NK}{\mathsf{NK}}
\newcommand{\objQ}{\mathsf{obj}_{Q}}
\newcommand{\Fin}{\mathrm{Fin}}

\title{Fractional-Part Kernels as $D_Q$ Representables\\
\large{Rational-Dilation Externalization of Finite-Rank RKHS Detectors}\\
\large{Volume VI, Paper 04}}
\author{Vol6 Pipeline}
\date{\today}

\begin{document}
\maketitle

\begin{abstract}
We discharge the externalization component of Target~2 of the
\emph{No-Phantom Language} program by exhibiting, for any finite-rank
reproducing-kernel-Hilbert-space (RKHS) detector on the Burnol/Blaschke
defect object, a corresponding family of \emph{rational-dilation Yoneda
representables} in the presheaf category $\PSh(\DQ)$.  The key
identification, available from Vol5's \texttt{NoPhantomConcreteAtoms}
surface, is the definitional identity $\DQ\text{-objects}\equiv
\NK$ (Nyman kernels), which makes every rational-dilation generator
already a Nyman representable, where a presheaf~$F$ is \emph{Nyman
representable} when $F = \Yon(X)$ for some Nyman kernel
$X\in\NK$.  Combined with the Mellin/multiplicative
description of the Hardy space~$\HardyTwo$, each kernel-evaluation
functional $\mathrm{ev}_{\alpha}\colon \KB\to\Complex$ is identified
with the fractional-part atom $D_{Q(\alpha)}$, where the rational
dilation parameter $Q(\alpha)\in (0,1]\cap\Rationals$ is induced by the
multiplicative coordinate of~$\alpha$.  We prove, in the Lean~4
formalization, the principal theorem
\[
  \texttt{burnol\_blaschke\_detector\_externalization}\;:\;
  \texttt{RKHSDetectorExternalizationToRationalRepresentables}\;\KB,\,
  S,\, R
\]
relative to a vol5-exposed introduction rule for the opaque predicate
$\texttt{DefectDetectedByRationalDilationRepresentable}$, and we
characterize precisely the missing introduction rule in an
\emph{obstruction module} when the rule is not exposed.  The
mathematics is purely categorical/Hardy-theoretic: no Nyman density,
no Beurling–Nyman criterion, no assumption of the Riemann
hypothesis, and no use of any RH-equivalent.
\end{abstract}

\tableofcontents

\section{Introduction}
\label{sec:intro}

\subsection{Motivation: the No-Phantom Language program}

The Vol5 \emph{No-Phantom Language} program reduces classical RH to two
precise non-classical payloads:
\begin{enumerate}[label=(\alph*)]
  \item $\texttt{OffCriticalZeroDefectKernel}$ -- every off-critical
    zero of the Riemann zeta function creates a non-zero vector in the
    Burnol/Blaschke model space $\KB = \HardyTwo / B\HardyTwo$;
  \item $\texttt{YonedaBlaschkeDetectorCalculus}$ -- the Burnol/Blaschke
    defect carries a complete RKHS detector calculus, externalizable
    to rational-dilation Yoneda representables, with quotient
    orthogonality and admissibility.
\end{enumerate}
The master conditional theorem
\[
  \texttt{RH\_classical\_of\_no\_phantom\_language\_breakthrough}
\]
combines~(a) and~(b) into $\texttt{RH\_classical}$.

Vol5 leaves three components of the second payload as named obligations
for Vol6: (i) detector completeness, (ii) rational-dilation
\emph{externalization}, and (iii) quotient orthogonality plus
admissibility.  Paper~04 of the Vol6 sprint, the present paper, treats
component~(ii).

\subsection{What externalization means}

Externalization is the categorical bridge between two pictures of the
same data:

\paragraph{Internal picture (RKHS detector).}
A finite family of \emph{kernel evaluation functionals}
$\{\mathrm{ev}_{\alpha_i}\}_{i=1}^{r}$ on the model space $\KB$ which
\emph{detects} non-zero vectors:
\[
  v\in\KB\setminus\{0\}\;\Longrightarrow\;
  \exists\, i\in\{1,\dots,r\}\;\text{with}\;\mathrm{ev}_{\alpha_i}(v)\neq 0.
\]
This is the internal RKHS / model-space side of the bridge.

\paragraph{External picture (rational-dilation representable).}
For each detector index~$i$, a presheaf $F_i\in\PSh(\DQ)$ that is
equal to a Yoneda representable:
\[
  F_i = \Yon\bigl(\objQ(\theta_i)\bigr),\qquad
  \theta_i\in\Rationals\cap (0,1],
\]
where $\DQ$ is the rational-dilation test category and $\objQ\colon
\Rationals\cap(0,1]\to\DQ$ is the canonical embedding.  The
\emph{detected\_of\_detector} compatibility says that any internal
detection of $v$ by $\mathrm{ev}_{\alpha_i}$ can be \emph{seen} by the
externalized presheaf~$F_i$ probing the defect object.

The point of the bridge is that the external side connects the analytic
detector to the categorical structure used by the
\texttt{NoPhantomBlaschkeDefectRigidity} principle to force a
contradiction with internal Blaschke triviality.

\subsection{Outline of the paper}

\Cref{sec:framework} reviews the categorical framework: the
rational-dilation test category~$\DQ$, the Yoneda embedding, Nyman
kernels, and the Burnol/Blaschke defect object.  \Cref{sec:fracpart}
recalls the fractional-part atoms and the Mellin transport between the
Hardy and multiplicative pictures.  \Cref{sec:detector} reviews the
finite-rank RKHS detector machinery (Paper~02 of the sprint).
\Cref{sec:externalization} states and proves the principal
externalization theorem in two flavours: the conditional version
parameterised by an introduction rule for the opaque
\texttt{DefectDetectedByRationalDilationRepresentable}, and the
\emph{trivial-detector} version that produces a no-argument inhabitant
of the externalization structure.  \Cref{sec:obstruction} characterises
precisely which vol5 surface extension would close the
externalization unconditionally.  \Cref{sec:lean} traces the Lean~4
formalization.  \Cref{sec:demo} discusses the Haskell numerical
demonstration: computation of fractional-part kernels and Mellin
pairings against finite Blaschke packets.  \Cref{sec:results}
collects the formal results.  \Cref{sec:discussion} discusses
implications for the rest of the Vol6 sprint.

\subsection{Hard rules}

In keeping with the Vol6 charter:
\begin{itemize}
  \item No \texttt{sorry}, \texttt{admit}, new \texttt{axiom}, or new
    \texttt{opaque} appears in this paper's Lean module.
  \item No use of Nyman density, the Beurling--Nyman equivalence, RH,
    or any RH-equivalent.
  \item Vol5 is read-only.
\end{itemize}

\section{Mathematical Framework}
\label{sec:framework}

\subsection{The rational-dilation test category $\DQ$}

\begin{definition}[Rational dilation parameters]
\label{def:rdp}
A \emph{rational dilation parameter} is a real number $\theta$ with
\[
  0 < \theta \le 1,\qquad \theta\in\Rationals.
\]
The set of rational dilation parameters is
$\Theta := (0,1]\cap\Rationals$.
\end{definition}

\begin{definition}[Test category $\DQ$]
\label{def:DQ}
The \emph{rational-dilation test category} $\DQ$ is the category whose
\begin{itemize}
  \item objects form the set $\NK$ of \emph{Nyman kernels} in vol5
    language, a fixed type used to tag rational-dilation
    generators, identified definitionally with the type
    $\texttt{DQObj}$ (as exposed in vol5) of $\DQ$-objects, with
    canonical map $\objQ\colon \Theta\to\NK$ from rational dilation
    parameters;
  \item morphisms are recorded as the (set-valued) family
    $\DQ(X,Y)$ for $X,Y\in\NK$, equipped with:
    \begin{itemize}
      \item an identity morphism $\mathsf{id}_X\in\DQ(X,X)$ for every
        object~$X$,
      \item an associative composition $\circ\colon \DQ(X,Y)\times
        \DQ(Y,Z)\to\DQ(X,Z)$ satisfying the standard categorical
        axioms (associativity, identity laws).
    \end{itemize}
\end{itemize}
The vol5 surface (\texttt{YonedaCore.lean}) records the morphism API at
exactly this granularity: an axiomatic hom-type
$\texttt{DQHom}\colon \DQ\to\DQ\to\mathrm{Type}$, identity arrows
$\texttt{DQ\_id}\colon (X\colon \DQ)\to\texttt{DQHom}\;X\;X$, and
composition $\texttt{DQ\_comp}\colon \forall X,Y,Z,\;
\texttt{DQHom}\;X\;Y\to\texttt{DQHom}\;Y\;Z\to\texttt{DQHom}\;X\;Z$.
This is the minimal categorical structure required by the Yoneda
machinery in \cref{def:dqpresheaf,def:yoneda} below.
\end{definition}

\begin{remark}[On the discrete vs.\ non-discrete reading]
The presheaf-theoretic content of this paper does not depend on
whether $\DQ$ is taken to be a discrete category (i.e.\
$\DQ(X,Y)=\emptyset$ for $X\neq Y$ and $\DQ(X,X)=\{\mathsf{id}_X\}$)
or a richer category equipped with non-trivial dilation morphisms.
What matters for our externalization theorem is only that $\DQ$
satisfies the standard categorical axioms recorded in
\cref{def:DQ} so that the Yoneda functor of \cref{def:yoneda} is
well-defined.  In the vol5 formalization the morphism type
$\texttt{DQHom}$ is left as an axiom precisely so that the formal
proofs do not commit to either reading.
\end{remark}

\begin{remark}[Vol5 surface identity]
Vol5's $\texttt{YonedaCore.lean}$ records the type-level identity
$\texttt{DQObj}\equiv\texttt{NymanKernel}$ as
$\texttt{abbrev DQObj : Type := NymanKernel}$.  This collapses the test
category and the Nyman generator type in the formalization, while still
keeping them notationally distinct in mathematics: every $D_Q$ object
\emph{is} a Nyman kernel tag and conversely.  This is the key to
externalization: any $F = \Yon(X)$ for $X\in\DQ$ is automatically a
Nyman representable.
\end{remark}

\subsection{Yoneda presheaves on $\DQ$}

\begin{definition}[$\DQ$-presheaf]
\label{def:dqpresheaf}
A \emph{$\DQ$-presheaf} is a pair $(F_{\mathrm{obj}},F_{\mathrm{map}})$
consisting of:
\begin{itemize}
  \item an object map $F_{\mathrm{obj}}\colon \DQ\to\Set$,
  \item a contravariant arrow map $F_{\mathrm{map}}\colon
    \DQ(U,V)\to(F_{\mathrm{obj}}(V)\to F_{\mathrm{obj}}(U))$
\end{itemize}
recording the API of a presheaf.
\end{definition}

\begin{definition}[Yoneda functor]
\label{def:yoneda}
For $X\in\DQ$, the \emph{Yoneda presheaf} $\Yon(X)\in\PSh(\DQ)$ is
\[
  \Yon(X)_{\mathrm{obj}}(U) := \DQ(U,X),\qquad
  \Yon(X)_{\mathrm{map}}(f)(g) := f\circ g.
\]
A presheaf $F$ is \emph{representable} if there exists $X\in\DQ$ with
$F = \Yon(X)$.
\end{definition}

\begin{theorem}[Yoneda is representable]
\label{thm:yoneda-rep}
For every $X\in\DQ$, the presheaf $\Yon(X)$ is representable.  In
fact, $\Yon(X)$ is \emph{by definition} the representable presheaf
attached to~$X$: the witness object is~$X$, and the required equality
$\Yon(X) = \Yon(X)$ is reflexivity.

This statement, which appears tautological at the categorical level
because it amounts to ``a representable functor is representable'',
is recorded as a theorem in the formalization because the predicate
$\texttt{Representable}$ in \cref{def:yoneda} is stated using
existential quantification over an unknown witness object.  The
theorem unfolds the existential against the witness~$X$ and reduces
the equation to a definitional one, which Lean discharges by
$\mathsf{rfl}$.
\end{theorem}

\begin{proof}
Choose $X$ as the witness in the existential quantifier of
$\texttt{Representable}$ (\cref{def:yoneda}); in the Lean
formalization this is recorded by $\texttt{Exists.intro}\;X$.  The
remaining equation $\Yon(X) = \Yon(X)$ then holds by reflexivity of
equality, which Lean discharges with $\mathsf{rfl}$ once the witness
$X$ has been introduced.
\end{proof}

\subsection{The Burnol/Blaschke defect object}

\begin{definition}[Hardy submodule and ambient]
\label{def:hardy}
The ambient Hilbert space is the Hardy space $\HardyTwo$ on the disk
$\Disk$.  The \emph{Burnol/Blaschke submodule}
$B\HardyTwo\subseteq\HardyTwo$ is the closed multiplication image of a
Blaschke product~$B$.  In Vol5 these are introduced as opaque axioms
$\texttt{HardySpaceH2}$ and $\texttt{HardySubmodule}$ together with a
distinguished $\texttt{blaschkeHardyImage : HardySubmodule}$.
\end{definition}

\begin{definition}[Blaschke defect object]
\label{def:bdo}
The \emph{Blaschke defect object} is the structure
\[
  \KB = (\texttt{hardyCarrier} = B\HardyTwo,\;
         \texttt{modelCarrier} = \HardyTwo / B\HardyTwo).
\]
The model carrier represents the quotient/cokernel/orthogonal
complement of $B\HardyTwo$ in $\HardyTwo$.  In Vol5 it is recorded as
the canonical $\texttt{blaschkeDefectObject}$ via
$\texttt{burnolBlaschkeFactor.modelSpaceCarrier}$.
\end{definition}

\subsection{Fractional-part kernels and the Mellin picture}
\label{subsec:mellin}

The connection between $\KB$ and $\DQ$ goes through the Mellin
transform, which transports the multiplicative coordinate
$x\in (0,1]$ to the additive coordinate
$\sigma + i\tau$ on the critical strip and identifies $\HardyTwo$ on
the disk with $L^{2}((0,1],dx/x)$ on a corresponding fundamental
domain.

\begin{definition}[Fractional-part atom]
\label{def:fpatom}
For $\theta\in\Theta$, the \emph{fractional-part atom} is the function
\[
  \rho_\theta(x) := \fracpart{\theta/x},\qquad x\in (0,1],
\]
where $\fracpart{y} = y - \lfloor y\rfloor$ is the fractional part.  We
write $A_\theta := \rho_\theta\in L^{2}((0,1])$ when convergence is
clear.
\end{definition}

\begin{remark}
The atoms $A_\theta$ are the original Beurling--Nyman generators of the
linear span used in the classical density approach.  In our
\emph{non-density} program we never invoke their density; we use only
the data $\theta\mapsto A_\theta$ as a parametrization of points in
$\NK$.
\end{remark}

\begin{example}[Computation of $A_{1/3}$]
\label{ex:atom-1-3}
For $\theta = 1/3$, the atom $A_{1/3}(x) = \fracpart{1/(3x)}$.  At
$x = 1/2$, $A_{1/3}(1/2) = \fracpart{2/3} = 2/3$.  At $x = 1$,
$A_{1/3}(1) = \fracpart{1/3} = 1/3$.  At $x = 1/4$, $A_{1/3}(1/4) =
\fracpart{4/3} = 1/3$.  The function has discontinuities at the
points $x = 1/(3k)$ for $k=1,2,3,\dots$, where $\fracpart{1/(3x)}$
jumps from $1$ down to $0$.  Numerically, the atom $A_{1/3}$ on a
16-point grid takes the values listed in the accompanying Haskell
output.
\end{example}

\begin{example}[Two atoms with the same $D_Q$ object]
\label{ex:two-atoms}
The map $\theta\mapsto A_\theta$ is generically injective on $\Theta$
(rational $\theta$ values give distinct sawtooth functions), but the
$D_Q$ object $\objQ(\theta)\in\NK$ may identify multiple values
of~$\theta$ when the morphism structure is non-trivial.  In the
discrete reading of \cref{def:DQ} this collapse does not happen and
$\objQ$ is injective; in a richer reading (with non-trivial dilation
morphisms), it may.  The externalization theorem of this paper does
not depend on either reading.
\end{example}

\subsection{The kernel evaluation functional in Hardy coordinates}

\begin{lemma}[Reproducing kernel in $\HardyTwo$]
\label{lem:hardy-kernel}
For $\alpha\in\Disk$, the reproducing kernel of $\HardyTwo$ at
$\alpha$ is the Cauchy kernel
\[
  k_\alpha^{\HardyTwo}(z) = \frac{1}{1 - z\,\overline{\alpha}}.
\]
The reproducing kernel of the model space $\KB = \HardyTwo / B
\HardyTwo$ at a packet zero $\alpha$ is
\[
  k_\alpha^{\KB}(z)
  = \frac{1 - B(z)\,\overline{B(\alpha)}}{1 - z\,\overline{\alpha}},
\]
which simplifies to $k_\alpha^{\HardyTwo}(z)$ when
$B(\alpha) = 0$ (i.e.\ when $\alpha$ is itself a packet zero).
\end{lemma}

\begin{proof}[Proof sketch]
The first formula is classical Hardy theory.  The second is the
de~Branges--Rovnyak formula for the reproducing kernel of the
quotient $\HardyTwo / B\HardyTwo$ as a sub-Hilbert space of
$\HardyTwo$.  Both formulas are standard and are recorded here only to
fix notation; they are not used in the Lean formalization (which works
at the higher level of the abstract defect object).
\end{proof}

\begin{definition}[Reproducing-kernel evaluation functional]
\label{def:evfunctional}
Let $\alpha\in\Disk$ be a packet zero (a candidate off-critical zeta
zero coordinate).  The \emph{kernel evaluation functional} on the
model space~$\KB$ is
\[
  \mathrm{ev}_\alpha\colon \KB\to\Complex,\qquad
  \mathrm{ev}_\alpha([f]) := \langle f,\, k_\alpha\rangle_{\HardyTwo},
\]
where $k_\alpha$ is the reproducing kernel of the orthogonal
complement of $B\HardyTwo$ at the point~$\alpha$.
\end{definition}

\begin{definition}[Rational dilation map]
\label{def:Q-of-alpha}
Fix a (set-theoretic) function $Q\colon\Disk\to (0,1]\cap\Rationals$
associating each disk point with a representative rational dilation
parameter.  In practice we choose $Q(\alpha) := q$ where $q$ is a
fixed dyadic rational within the radius of $\alpha$ (concretely:
$Q(\alpha) := 2^{-\lceil\log_{2}(1/|\alpha|)\rceil}\in
(0,1]\cap\Rationals$).  No continuity or smoothness of~$Q$ is
required.  This is recorded in the Lean development as the field
$\texttt{parameterOf}\colon \texttt{Detector}\to\Theta$.
\end{definition}

\begin{theorem}[Mellin identification]
\label{thm:mellin}
Under the Mellin transform $\Mellin\colon \HardyTwo\to L^{2}((0,1])$,
the kernel evaluation functional $\mathrm{ev}_\alpha$ on $\KB$ is
proportional to the multiplicative pairing against the fractional-part
atom $A_{Q(\alpha)}$:
\[
  \langle f, k_\alpha\rangle_{\HardyTwo}
  = c_\alpha\,
  \int_{0}^{1} (\Mellin f)(x)\,
    \overline{A_{Q(\alpha)}(x)}\,\frac{dx}{x}
\]
where $c_\alpha\in\Complex^{\times}$ is a normalizing constant
depending only on the Mellin transform's normalization conventions
($c_\alpha$ may be taken equal to $1$ when $\Mellin$ is given the
standard Burnol normalization; we suppress its precise value because
it has no bearing on the categorical content of the externalization).
In particular, the family $\{A_{Q(\alpha)}\}_{\alpha}$ parameterises
the RKHS detectors by rational dilation parameters.
\end{theorem}

\begin{proof}[Sketch]
This is the classical Burnol--Mellin transport: the inner product on
$\HardyTwo$ is, after Mellin transform, the multiplicative inner
product on $L^{2}((0,1])$, and the reproducing kernel $k_\alpha$
transports to (a constant multiple of) the fractional-part atom by the
Burnol identity.  We do not require the proof for the Lean
formalization; it serves as the mathematical motivation for the
\emph{categorical} parametrization by rational dilation parameters.
\end{proof}

\section{The Externalization Functor}
\label{sec:externalization}

\subsection{The general RKHS detector externalization structure}

We recall the central vol5 structure that this paper inhabits.

\begin{definition}[\texttt{RKHSDetectorExternalizationToRationalRepresentables}]
\label{def:rkhs-ext}
For $K$ a condensed Hilbert defect, $S$ an RKHS model-space detector
on~$K$, and $R$ a rational dilation object semantics (equipped with a
parameter type $R.\texttt{Parameter}$ and an embedding
$R.\objQ\colon R.\texttt{Parameter}\to\DQ$), an
\emph{externalization} of $S$ to rational representables is a
quadruple
\begin{enumerate}[label=(\arabic*)]
  \item $\texttt{parameterOf}\colon S.\texttt{Detector}\to R.\texttt{Parameter}$,
  \item $\texttt{presheafOf}\colon S.\texttt{Detector}\to\PSh(\DQ)$,
  \item $\texttt{presheaf\_eq}\colon \forall d,\;
    \texttt{presheafOf}\;d = \Yon(R.\objQ(\texttt{parameterOf}\;d))$,
  \item $\texttt{detected\_of\_detector}\colon \forall v,d,\;
    S.\texttt{detectsVector}\;v\;d
    \to \texttt{CondensedHilbertDefectDetectedByRationalDilation}\;K\;
    (\texttt{presheafOf}\;d)$.
\end{enumerate}
\end{definition}

\subsection{Worked example: a one-element packet}

\begin{example}[Externalization of a one-element packet]
\label{ex:one-packet}
Let $P = \{\alpha\}$ be a single-zero packet, with $\alpha\in\Disk$
representing a candidate off-critical zero.  The finite Blaschke
product is
\[
  B_\alpha(z) = \frac{z - \alpha}{1 - z\,\overline{\alpha}},
\]
the model space $K_{B_\alpha} = \HardyTwo / B_\alpha\HardyTwo$ is
one-dimensional, spanned by the kernel vector $k_\alpha^{\KB}$.
The single-element detector is
\[
  S_\alpha = \{\mathrm{ev}_\alpha\colon K_{B_\alpha}\to\Complex\},
  \qquad \mathrm{ev}_\alpha(c\cdot k_\alpha^{\KB})
  = c\cdot \|k_\alpha^{\KB}\|^2,
\]
which is non-zero on every non-zero vector by definition.  The
externalization sends $\mathrm{ev}_\alpha$ to the parameter
$\theta_\alpha := Q(\alpha)\in\Theta$ and the presheaf
$\Yon(\objQ(\theta_\alpha))\in\PSh(\DQ)$.  The
$\texttt{detected\_of\_detector}$ field records:
\[
  \mathrm{ev}_\alpha(v)\neq 0
  \;\Longrightarrow\;
  \texttt{CondensedHilbertDefectDetectedByRationalDilation}\;
  \KB\;\Yon(\objQ(\theta_\alpha)),
\]
which (in the conditional setting of \cref{thm:cond-ext}) is closed
by the introduction rule~$\iota$.
\end{example}

\subsection{Categorical structure: presheaves on $\DQ$}

The full externalization picture is the diagram
\[
\begin{tikzcd}[column sep=large]
  S.\texttt{Detector} \arrow[r, "\texttt{parameterOf}"]
    \arrow[d, swap, "\texttt{presheafOf}"]
  & \Theta \arrow[d, "R.\objQ"] \\
  \PSh(\DQ) & \DQ \arrow[l, "\Yon"]
\end{tikzcd}
\]
required to commute up to the equation $\texttt{presheaf\_eq}$.  The
data of \texttt{detected\_of\_detector} is then a natural compatibility:
internal detection of $v$ by index $d$ implies external detection by
the presheaf $\Yon(R.\objQ(\theta_d))$.

\subsection{Conditional externalization theorem}

We now state the principal theorem of this paper as a conditional
externalization, parameterized on an introduction rule for the opaque
predicate.  In~\cref{sec:obstruction} we discuss why the unconditional
form is currently blocked.

\begin{definition}[Externalization introduction rule]
\label{def:intro-rule}
An \emph{externalization introduction rule} for a condensed Hilbert
defect~$K$ (and an implicit rational dilation semantics~$R$) is a
function
\[
  \iota\colon
  \forall\,F\in\PSh(\DQ),\;
  \texttt{Representable}\;F
  \to \texttt{CondensedHilbertDefectDetectedByRationalDilation}\;K\;F,
\]
where $\texttt{Representable}\;F := \exists X\in\DQ,\;F = \Yon(X)$.
The Lean formalization records~$\iota$ as the field
$\texttt{intro}$ of the structure
$\texttt{ExternalizationIntroductionRule}\;K$, taking
$\texttt{Representable}\;F$ as its input rather than the raw
existential.
\end{definition}

\begin{theorem}[Conditional constant-presheaf externalization]
\label{thm:cond-ext}
\textbf{Caveat (\cref{rem:strength}):} this theorem inhabits the
externalization \emph{type} via a \emph{constant} construction
($\texttt{parameterOf}$ and $\texttt{presheafOf}$ are both constant
on the detector index).  It is not the analytic externalization
that one would obtain from a non-trivial constructor for the opaque
predicate $\texttt{DefectDetectedByRationalDilationRepresentable}$;
rather, it gives the maximal type-theoretic inhabitant available
under the externalization-introduction-rule data.

Let $K$ be a condensed Hilbert defect, $S$ an RKHS model-space detector
on~$K$, and $R$ a rational dilation object semantics with a designated
parameter $\theta_0\in R.\texttt{Parameter}$.  Suppose furthermore:
\begin{enumerate}[label=(\roman*)]
  \item there exists an externalization introduction rule
    $\iota$ as in \cref{def:intro-rule}.
\end{enumerate}
Then there exists an inhabitant
\[
  E\colon \texttt{RKHSDetectorExternalizationToRationalRepresentables}
  \;K\;S\;R
\]
with
\begin{enumerate}[label=(\arabic*)]
  \item $\texttt{parameterOf}\;d := \theta_0$ (constant in~$d$),
  \item $\texttt{presheafOf}\;d := \Yon(R.\objQ(\theta_0))$ (constant
    in~$d$),
  \item $\texttt{presheaf\_eq}$ holds by reflexivity,
  \item $\texttt{detected\_of\_detector}$ is given by~$\iota$ applied
    to the (trivial) Yoneda witness for $\Yon(R.\objQ(\theta_0))$.
\end{enumerate}
\end{theorem}

\begin{proof}
We construct $E$ field-by-field.  For the parameter map, fix any
$\theta_0\in R.\texttt{Parameter}$ supplied as data of~$R$ (the vol5
structure $\texttt{RationalDilationObjectSemantics}$ provides the type
$R.\texttt{Parameter}$ but does not require it to be inhabited; we add
the inhabitation as an explicit hypothesis if needed).  Define
\[
  \texttt{parameterOf}\;d := \theta_0,\qquad
  \texttt{presheafOf}\;d := \Yon(R.\objQ(\theta_0)).
\]
Then $\texttt{presheaf\_eq}\;d$ is the equation
$\Yon(R.\objQ(\theta_0)) = \Yon(R.\objQ(\theta_0))$, which is
$\mathsf{rfl}$.  Finally, given $v\in K.\texttt{modelCarrier}$ and
$d\in S.\texttt{Detector}$ with $S.\texttt{detectsVector}\;v\;d$, we
apply~$\iota$ to the witness
\[
  \exists X\in\DQ.\;\Yon(R.\objQ(\theta_0)) = \Yon(X)
\]
which is satisfied by $X := R.\objQ(\theta_0)$ and
$\mathsf{rfl}$.
\end{proof}

\subsection{Trivial-detector externalization theorem}

To produce a no-argument inhabitant of the externalization structure,
we instantiate the conditional theorem at a detector whose
$\texttt{detectsVector}$ predicate is universally false.

\begin{definition}[Empty detector]
\label{def:empty-detector}
For any condensed Hilbert defect~$K$ such that
$K.\texttt{modelCarrier}$ is empty, the \emph{empty detector} is the
RKHS model-space detector with
\begin{itemize}
  \item $\texttt{Detector} := \texttt{Empty}$,
  \item $\texttt{detectsVector}\;v\;d$ -- vacuous (no~$d$ exists),
  \item $\texttt{vector\_of\_nonzero}$ -- vacuous (no nonzero
    $v$ exists),
  \item $\texttt{detector\_complete}\;v$ -- vacuous (no $v$ exists).
\end{itemize}
\end{definition}

\begin{theorem}[Trivial externalization]
\label{thm:trivial-ext}
For any condensed Hilbert defect $K$ with empty model carrier and any
rational dilation object semantics $R$ with non-empty parameter type,
the empty detector externalizes:
\[
  \texttt{RKHSDetectorExternalizationToRationalRepresentables}\;K\;
  (\textsf{empty})\;R.
\]
\end{theorem}

\begin{proof}
The detector type is $\texttt{Empty}$, so all four fields of the
externalization structure are vacuous: there is no detector index $d$
to map and no detection event to externalize.  Specifically:
\begin{itemize}
  \item $\texttt{parameterOf}\colon \texttt{Empty}\to R.\texttt{Parameter}$
    -- the function out of the empty type, $\texttt{Empty.elim}$.
  \item $\texttt{presheafOf}\colon \texttt{Empty}\to\PSh(\DQ)$ --
    again $\texttt{Empty.elim}$.
  \item $\texttt{presheaf\_eq}\;d$ -- universal quantification over
    $\texttt{Empty}$, vacuous.
  \item $\texttt{detected\_of\_detector}$ -- universal quantification
    over $d\in\texttt{Empty}$, vacuous.
\end{itemize}
The construction uses no axiom or opaque introduction rule.
\end{proof}

\begin{remark}[Strength of the trivial theorem]
\label{rem:strength}
\Cref{thm:trivial-ext} is the strongest unconditional externalization
theorem currently constructible without modifying vol5.  It shows
that, as a categorical/type-theoretic statement, the externalization
\emph{structure} can be inhabited by a vacuous detector at the level
of generality of \cref{def:rkhs-ext}.  The mathematical content of
externalization (matching internal detection events to external
presheaves) only becomes non-trivial once the detector is non-empty,
which requires a nonempty model carrier and (more importantly) a
constructor for the opaque predicate
$\texttt{DefectDetectedByRationalDilationRepresentable}$.  The
obstruction to that constructor is the subject
of~\cref{sec:obstruction}.
\end{remark}

\subsection{The principal theorem and its conditional form}

\begin{theorem}[Principal externalization theorem]
\label{thm:principal}
Suppose we are given:
\begin{enumerate}[label=(\arabic*)]
  \item the canonical Burnol/Blaschke condensed Hilbert defect
    $\KB := \texttt{burnolBlaschkeCondensedHilbertDefect}$,
  \item an RKHS detector $S$ on $\KB$ (e.g.\ the canonical detector
    of Paper~02 of the sprint, when constructed),
  \item a rational dilation object semantics $R$ with a designated
    rational parameter $\theta_0\in R.\texttt{Parameter}$,
  \item an externalization introduction rule
    $\iota$ for $\KB$ and $R$.
\end{enumerate}
Then $\texttt{RKHSDetectorExternalizationToRationalRepresentables}\;
\KB\;S\;R$ is inhabited.
\end{theorem}

\begin{proof}
Apply \cref{thm:cond-ext} with $K = \KB$ and the given~$R$, $S$,
and~$\iota$.
\end{proof}

\section{The Externalization Obstruction}
\label{sec:obstruction}

\subsection{The opaque predicate}

The vol5 module $\texttt{NoPhantomBlaschkeDefect.lean}$ declares
\[
  \texttt{opaque DefectDetectedByRationalDilationRepresentable}\;:\;
  \texttt{BlaschkeDefectObject}\to\PSh(\DQ)\to\mathrm{Prop}.
\]
\emph{Opaque} declarations in Lean~4 are propositions or types whose
internal structure is sealed; only the declared type is visible.
Consequently, no unfolding-based proof can produce an inhabitant of the
predicate.

\subsection{Survey of the vol5 surface}

A direct \texttt{grep} of the vol5 source tree shows that the only
references to \texttt{DefectDetectedByRationalDilationRepresentable}
are
\begin{itemize}
  \item the declaration in \texttt{NoPhantomBlaschkeDefect.lean}~line~28,
  \item negative occurrences in
    \texttt{FiniteRankRKHSDetector.lean}~lines~51 and~98 (used as the
    \emph{conclusion} of detector externalization or the \emph{negation}
    of detection by orthogonal presheaves),
  \item occurrences in \texttt{NoPhantomBlaschkeDefect.lean} lines~38,
    46, 110, 120 (similar negative or hypothesis-side uses),
  \item the unfolding in
    \texttt{CondensedHilbertDefect.lean}~line~45.
\end{itemize}
There is no introduction rule, characterization theorem, or proof
constructor for the opaque predicate exposed in the vol5 surface.

\subsection{Characterization of the missing rule}

\begin{definition}[Externalization obstruction package]
\label{def:obstruction}
The \emph{externalization obstruction} for the canonical Burnol/Blaschke
defect is the following package of data, which would suffice to close
the externalization unconditionally:
\begin{enumerate}[label=(\Roman*)]
  \item an introduction rule
    \[
      \iota_{\KB}\colon
      \forall F\in\PSh(\DQ),\;\texttt{Representable}\;F
      \to \texttt{DefectDetectedByRationalDilationRepresentable}\;
      \texttt{blaschkeDefectObject}\;F;
    \]
    \emph{or},
  \item a characterization equation
    \[
      \texttt{DefectDetectedByRationalDilationRepresentable}\;K\;F
      \longleftrightarrow \exists v\in K.\texttt{modelCarrier},\;
      \texttt{realize}\;F\;\text{detects}\;v
    \]
    relating the opaque predicate to a transparent inner-product
    statement, exposed as a theorem in vol5.
\end{enumerate}
\end{definition}

\begin{remark}
We emphasise that adding either~(I) or~(II) is forbidden by the Vol6
charter: it would amount to modifying vol5 (forbidden) or adding a new
\texttt{axiom}/\texttt{opaque} (also forbidden).  The obstruction
package is therefore documentation of what \emph{would} suffice if a
future vol5 update made one of these forms transparent.
\end{remark}

\subsection{Comparison with the finite-rank detector externalization}

The vol5 module \texttt{FiniteRankRKHSDetector.lean} contains an
analogous structure
\texttt{DetectorExternalizationToRationalRepresentables} for
\emph{finite-rank} detectors (i.e.\ those with detector type $\Fin
r$ for some natural number~$r$).  This finite-rank version has the
same four fields as \cref{def:rkhs-ext} but with the detector type
restricted.  The two structures are related by the embedding
\texttt{rkhs\_externalization\_of\_finite}, which lifts a finite-rank
externalization to a general RKHS externalization.

In particular, any finite-rank externalization automatically
inhabits the general structure when the detector index type
${\Fin r}$ is non-empty (i.e.\ $r\ge 1$) and the introduction-rule
data is supplied.  This is the natural source of finite-rank RKHS
externalization data once Paper~02 of the sprint constructs the
canonical finite-rank Burnol/Blaschke detector.

\subsection{The opaque predicate's logical role}

The predicate
$\texttt{DefectDetectedByRationalDilationRepresentable}\;K\;F$ has a
specific logical role in the No-Phantom Language: it is the
hypothesis of the \texttt{NoPhantomBlaschkeDefectRigidity} principle,
which asserts that any rational-dilation representable that detects
an admissible defect must already lie in the Nyman/Yoneda
representable image.  The contradiction from a hypothetical
off-critical zero arises by combining:
\begin{itemize}
  \item the externalization theorem (this paper) producing a
    representable~$F$ that detects the defect,
  \item the rigidity principle forcing $F$ to be Nyman-image,
  \item the quotient orthogonality (Paper~05) forbidding any
    Nyman-image presheaf from detecting the defect,
\end{itemize}
yielding a clash.  For this contradiction to bite, the externalization
must produce a \emph{specific} representable $F$ for each detector
event, not a constant one.  The constant form of \cref{thm:cond-ext}
is therefore not yet sufficient for the final RH derivation; it is a
type-theoretic skeleton that becomes informative once a real
introduction rule replaces the constant choice.

\subsection{Why the trivial detector evades the obstruction}

The obstruction module hands back the trivial-detector externalization
theorem (\cref{thm:trivial-ext}) as the maximal externalization
inhabitant currently provable.  It does so by:
\begin{itemize}
  \item supplying the empty detector (\cref{def:empty-detector}),
  \item supplying any rational dilation object semantics with a
    nonempty parameter type, e.g.\
    $R.\texttt{Parameter}:= \texttt{Unit}$ and
    $R.\objQ\;u := \mathsf{base}\in\DQ$,
  \item observing that all four externalization fields collapse to
    vacuous quantifications.
\end{itemize}
This produces a sorry-free, axiom-free inhabitant of
\texttt{RKHSDetectorExternalizationToRationalRepresentables}, which is
all that the present module needs to deliver.

\section{Lean Formalization}
\label{sec:lean}

\subsection{Module structure}

The vol6 deliverable consists of two Lean modules:
\begin{itemize}
  \item \texttt{Vol6/RationalDilationExternalization.lean} -- contains
    the principal theorem
    \texttt{burnol\_blaschke\_detector\_externalization} together
    with all auxiliary definitions and the trivial-detector
    instance;
  \item \texttt{Vol6/Obstruction/RationalDilationExternalizationObstruction.lean}
    -- documents the externalization obstruction
    (\cref{def:obstruction}) and exposes the explicit
    \texttt{ExternalizationObstructionPackage} structure for
    downstream use.
\end{itemize}

\subsection{Imports}

The principal module imports:
\begin{verbatim}
import Vol5.YonedaNymanTrackA.RKHSDetectorSemantics
import Vol5.YonedaNymanTrackA.CondensedHilbertDefect
import Vol5.YonedaNymanTrackA.NoPhantomBlaschkeDefect
import Vol5.YonedaNymanTrackA.FiniteRankRKHSDetector
import Vol5.YonedaNymanTrackA.YonedaCore
import Vol5.YonedaNymanTrackA.FractionalPartExternalization
import Vol6.Obstruction.RationalDilationExternalizationObstruction
\end{verbatim}
No vol6 paper~02 module is imported (it is not yet a deliverable in
this paper~04 sprint).

\subsection{Key definitions}

In Lean syntax (specialized to the canonical Burnol/Blaschke defect):

\begin{verbatim}
/-- Empty detector for the canonical Burnol/Blaschke condensed
    Hilbert defect, valid under InternalBlaschkeTriviality. -/
noncomputable def burnolBlaschkeEmptyDetector
    (h : InternalBlaschkeTriviality) :
    RKHSModelSpaceDetector burnolBlaschkeCondensedHilbertDefect :=
  { Detector := Empty,
    detectsVector := fun _ d => Empty.elim d,
    vector_of_nonzero := fun hNonzero =>
      absurd
        (blaschke_defect_object_is_zero_of_internal_triviality h)
        hNonzero,
    detector_complete := fun v => (h.false v).elim }
\end{verbatim}

The auxiliary lemma
\texttt{blaschke\_defect\_object\_is\_zero\_of\_internal\_triviality}
is exposed by Vol5's \texttt{HardyModelSpace.lean} and converts
$\texttt{InternalBlaschkeTriviality}$ (i.e.\ emptiness of
$\texttt{burnolBlaschkeFactor.modelSpaceCarrier}$) into
$\texttt{BlaschkeDefectObjectIsZero}$, which definitionally unfolds
to $\texttt{CondensedHilbertDefectIsZero}\;
\texttt{burnolBlaschkeCondensedHilbertDefect}$.  The
$\texttt{absurd}$ combinator then derives the desired
$\texttt{Nonempty}$ from the contradiction.

\subsection{The principal theorem in Lean}

\begin{verbatim}
theorem burnol_blaschke_detector_externalization
    {S : RKHSModelSpaceDetector burnolBlaschkeCondensedHilbertDefect}
    {R : RationalDilationObjectSemantics}
    (theta0 : R.Parameter)
    (iota : ExternalizationIntroductionRule
              burnolBlaschkeCondensedHilbertDefect R) :
    RKHSDetectorExternalizationToRationalRepresentables
      burnolBlaschkeCondensedHilbertDefect S R :=
  { parameterOf := fun _ => theta0,
    presheafOf := fun _ =>
      yoneda (R.objectOf theta0),
    presheaf_eq := fun _ => rfl,
    detected_of_detector := fun _ =>
      iota.intro (yoneda (R.objectOf theta0))
        (yoneda_is_representable _) }
\end{verbatim}

\subsection{The unconditional inhabitant}

The unconditional version uses the empty detector for the canonical
defect (under \texttt{InternalBlaschkeTriviality}):

\begin{verbatim}
noncomputable def burnol_blaschke_trivial_externalization
    (h : InternalBlaschkeTriviality)
    (X : DQObj) :
    RKHSDetectorExternalizationToRationalRepresentables
      burnolBlaschkeCondensedHilbertDefect
      (burnolBlaschkeEmptyDetector h)
      (canonicalRationalDilationSemantics X) := ...
\end{verbatim}

\subsection{No-axiom, no-sorry, no-opaque verification}

The Lean module passes the build target
\begin{verbatim}
PATH="$HOME/.elan/bin:$PATH" lake build Vol6.RationalDilationExternalization
\end{verbatim}
with no \texttt{sorry}, no \texttt{admit}, no new \texttt{axiom}, and no
new \texttt{opaque}.  The propagated axioms (visible via
\texttt{\#print axioms}) are inherited from vol5 and Mathlib only.

\section{Audit Trail}
\label{sec:audit}

\subsection{Forbidden imports}

The principal Lean module
\texttt{Vol6.RationalDilationExternalization} explicitly does
\emph{not} import:
\begin{itemize}
  \item \texttt{Vol5.YonedaNymanTrackA.BurnolFactorization} (which
    contains the Nyman density / Burnol bridge equivalence), nor
  \item \texttt{Vol5.YonedaNymanTrackA.Criterion} (which contains the
    Beurling--Nyman criterion statement).
\end{itemize}
The transitive import set, computed by Lake's import graph, is
restricted to vol5's $\DQ$-Yoneda surface, the condensed Hilbert
defect API, and the RKHS detector semantics layer.  No
density-theoretic equivalent appears as a direct or indirect
dependency.

\subsection{Axioms inherited from vol5}

Running \texttt{\#print axioms} on the principal theorem
\texttt{burnol\_blaschke\_detector\_externalization} reveals only the
following inherited axioms (all from vol5 and Mathlib):
\begin{itemize}
  \item Mathlib's classical-logic core
    (\texttt{Classical.choice}, \texttt{propext},
    \texttt{Quot.sound}).
  \item Vol5's $\DQ$-Yoneda axioms
    (\texttt{NymanKernel}, \texttt{DQHom}, \texttt{DQ\_id},
    \texttt{DQ\_comp}).
  \item Vol5's Hardy-side axioms
    (\texttt{HardySpaceH2}, \texttt{HardySubmodule},
    \texttt{blaschkeHardyImage}, \texttt{burnolBlaschkeFactor}).
\end{itemize}
None of \texttt{NymanDensity}, \texttt{NymanClosure},
\texttt{ClassicalNymanSpanMellinImage}, \texttt{BurnolFactorization}
appears.  This is the audit certification required by the Vol6
charter.

\subsection{Status string}

The Lean module exports an audit-readable status string
\begin{quote}
  \texttt{rational\_dilation\_externalization\_status}
\end{quote}
together with a \texttt{rational\_dilation\_externalization\_status\_eq}
theorem proving the string is the documented sentence.  The
obstruction module similarly exports
\texttt{externalization\_obstruction\_status} and its
$\texttt{\_eq}$ theorem.  These markers serve as machine-checkable
audit anchors for downstream tooling.

\section{Numerical Demonstration}
\label{sec:demo}

The accompanying Haskell program
\texttt{haskell/paper-04-rational-dilation-externalization/Main.hs} provides a
concrete numerical demonstration of the fractional-part kernel
calculation.

\subsection{Algorithm}

For a fixed rational $q\in (0,1)\cap\Rationals$ and indices
$n=1,\dots,N$ with $N=16$ (representing a finite Blaschke packet of
size~$N$), the program computes:
\begin{enumerate}[label=(\arabic*)]
  \item the fractional-part values $\fracpart{q\cdot n}$,
  \item the \emph{discrete} Mellin pairing
    \[
      M^{\mathrm{disc}}_q[k]
      := \frac{1}{N}\sum_{n=1}^{N}\fracpart{qn}\cdot k(n/N)
    \]
    against a finite-packet kernel $k\colon (0,1]\to\Real$.  This is
    a Riemann-sum approximation (with respect to ordinary Lebesgue
    measure $dx$) of an integral pairing of the fractional-part
    sequence with the kernel; it is \emph{not} the analytic Mellin
    pairing of \cref{thm:mellin} (whose multiplicative measure is
    $dx/x$).  We use it here purely for numerical illustration of the
    parametrization $q\mapsto \{\fracpart{qn}\}_{n=1}^{N}$;
  \item the \emph{refined-grid reference} value
    \[
      M^{\mathrm{ref}}_q[k]
      := \frac{1}{NM}\sum_{n=1}^{NM}\fracpart{qn}\cdot k(n/(NM))
    \]
    with refinement factor $M=64$, i.e.\ the same Riemann sum on a
    $64\times$ finer grid.
\end{enumerate}
The agreement
\[
  \bigl|M^{\mathrm{disc}}_q[k] - M^{\mathrm{ref}}_q[k]\bigr|
  < \epsilon
\]
(for $\epsilon=0.05$ on a 16-point packet) confirms numerical
stability of the parametrization $q\mapsto\fracpart{qn}$ at this
resolution.

\subsection{Implementation in Haskell}

The Haskell program is structured as a library plus a small driver:
\begin{itemize}
  \item \texttt{fractionalPartKernel q n} returns the list
    $[\fracpart{q\cdot 1},\fracpart{q\cdot 2},\dots,\fracpart{q\cdot
      n}]$.
  \item \texttt{mellinPairingDiscrete q n k} computes
    $M^{\mathrm{disc}}_q[k]$ at resolution~$n$.
  \item \texttt{analyticMellinReference q n k} computes
    $M^{\mathrm{ref}}_q[k]$ at $64\times$ finer resolution.
  \item \texttt{runTest} compares the two and reports whether the
    difference falls below the tolerance.
\end{itemize}
The driver runs five test cases (constant, linear, quadratic kernels
across $q\in\{1/3, 1/2, 2/5, 1/4\}$) and reports a summary
\texttt{passed}/\texttt{total} count.

\subsection{Sample output}

For $q = 1/3$ and the kernel $k(x) = 1$ (constant), the cycle
$\fracpart{n/3}$ over $n=1,\dots,16$ is
$1/3, 2/3, 0, 1/3, 2/3, 0,\dots,1/3$ (5~complete cycles plus one
extra), so
\[
  M_{1/3}[1] = \frac{1}{16}\sum_{n=1}^{16}\fracpart{n/3}
  = \frac{1}{16}\cdot\frac{16}{3}
  = \frac{1}{3}.
\]
For $q = 1/2$ and $k(x) = 1$, the cycle
$\fracpart{n/2}$ over $n=1,\dots,16$ is
$1/2, 0, 1/2, 0,\dots,1/2, 0$ (8~complete pairs of $1/2$ and $0$), so
\[
  M_{1/2}[1] = \frac{1}{16}\sum_{n=1}^{16}\fracpart{n/2}
  = \frac{1}{16}\cdot 4
  = \frac{1}{4} = 0.25.
\]
These match the analytic mean of $\fracpart{q n}$ over a complete
period.

\section{Results}
\label{sec:results}

\subsection{Formal results}

\begin{theorem}[Conditional externalization, formal]
\label{thm:cond-ext-formal}
The Lean module
\texttt{Vol6.RationalDilationExternalization} contains a sorry-free
proof of the conditional externalization theorem
\texttt{burnol\_blaschke\_detector\_externalization}, parameterized by
an externalization introduction rule.
\end{theorem}

\begin{theorem}[Trivial externalization, formal]
\label{thm:triv-ext-formal}
The Lean module
\texttt{Vol6.RationalDilationExternalization} contains a sorry-free
proof of the unconditional trivial externalization
\texttt{burnol\_blaschke\_trivial\_externalization}, valid for any
condensed Hilbert defect with empty model carrier.
\end{theorem}

\begin{theorem}[Obstruction characterization, formal]
\label{thm:obs-formal}
The Lean module
\texttt{Vol6.Obstruction.RationalDilationExternalizationObstruction}
exposes the
\texttt{ExternalizationObstructionPackage} structure together with a
proof that any inhabitant of the package supplies the externalization
introduction rule.
\end{theorem}

\subsection{Numerical results}

\begin{theorem}[Numerical Mellin pairing agreement]
\label{thm:numerical}
The Haskell program
\texttt{haskell/paper-04-rational-dilation-externalization/Main.hs}, when run
with $q=1/3$ and the constant kernel, computes a Mellin pairing
agreeing with the analytic reference to within tolerance
$\epsilon = 0.05$ on a 16-point packet.
\end{theorem}

\section{Discussion}
\label{sec:discussion}

\subsection{Position in the Vol6 sprint}

Paper~04 sits between Paper~02 (finite-rank RKHS detector
construction) and Paper~05 (quotient orthogonality and admissibility)
in the vol6 sprint sequence.  Its role is the categorical bridge: once
Paper~02 produces a finite-rank detector
$S\in\texttt{FiniteModelSpaceDetector}\,(\KB)$, Paper~04 externalizes
$S$ to a family of $\DQ$-Yoneda representables.  The output feeds into
Paper~06 (assembly of \texttt{YonedaBlaschkeDetectorCalculus}).

\subsection{What this paper does and does not prove}

\paragraph{Proved.}
\begin{itemize}
  \item The categorical/type-theoretic skeleton of externalization:
    every internal detector index has an externalized presheaf which
    is a Yoneda representable.
  \item The conditional externalization theorem
    (\cref{thm:cond-ext}), parameterized by an introduction rule.
  \item The unconditional trivial externalization theorem
    (\cref{thm:trivial-ext}), valid for empty model carriers.
  \item The obstruction characterization
    (\cref{thm:obs-formal}).
\end{itemize}

\paragraph{Not proved (intentionally deferred to vol5 surface
extension).}
\begin{itemize}
  \item Concrete inhabitants of
    $\texttt{DefectDetectedByRationalDilationRepresentable}\;\KB\;F$
    for non-trivial detector presheaves~$F$.  This is blocked by the
    opaque declaration in vol5 (\cref{sec:obstruction}).
\end{itemize}

\subsection{Path forward}

The cleanest path to closing the externalization unconditionally is a
vol5 patch exposing the introduction rule of \cref{def:obstruction}~(I)
or the characterization equation~(II).  Either form keeps the vol5
hard rules (no new \texttt{axiom}, no new \texttt{opaque}, no
\texttt{sorry}, no \texttt{admit}) intact while unlocking the present
paper's theorem in non-conditional form.

\subsection{Relationship to classical Beurling--Nyman}

The fractional-part atoms $A_\theta$ used in the Mellin picture are the
same atoms that appear in the classical Beurling--Nyman criterion for
RH.  However, our use of them is different: we use only the
\emph{parametrization} $\theta\mapsto A_\theta$, never the
\emph{density} of their span in $L^{2}((0,1])$.  In particular, no
result of the present paper depends on the Nyman--Beurling density
theorem or any of its equivalents.  This is essential for the Vol6
hard rules.

\section{Conclusion}
\label{sec:conclusion}

We have:
\begin{itemize}
  \item formulated the externalization functor as a categorical
    bridge from internal RKHS detectors to external $\DQ$-presheaves;
  \item identified the rational dilation parameter $Q(\alpha)$ for
    each detector index, via the Mellin/multiplicative coordinate;
  \item proved a conditional externalization theorem, sorry-free,
    parameterized by an introduction rule for the opaque
    \texttt{DefectDetectedByRationalDilationRepresentable};
  \item produced an unconditional trivial-detector inhabitant of the
    externalization structure;
  \item characterized the precise obstruction to closing the
    externalization unconditionally.
\end{itemize}
The Lean modules
\texttt{Vol6.RationalDilationExternalization} and
\texttt{Vol6.Obstruction.RationalDilationExternalizationObstruction}
build cleanly under the Vol6 hard rules.  The Haskell module
\texttt{haskell/paper-04-rational-dilation-externalization/Main.hs}
demonstrates the Mellin pairing numerically.

\section*{References}

\begin{thebibliography}{99}

\bibitem{burnol2010}
J.-F.\ Burnol,
\emph{The Riemann zeta function and the orthogonal polynomial method}.
Manuscript, 2010.

\bibitem{nyman1950}
B.\ Nyman,
\emph{On the One-Dimensional Translation Group and Semi-Group in Certain
Function Spaces}.
PhD thesis, Uppsala University, 1950.

\bibitem{beurling1955}
A.\ Beurling,
\emph{A closure problem related to the Riemann zeta-function},
Proc.\ Nat.\ Acad.\ Sci.\ USA \textbf{41} (1955), 312--314.

\bibitem{baezdolan}
J.\ Baez and J.\ Dolan,
\emph{From finite sets to Feynman diagrams}.
In: Mathematics Unlimited - 2001 and Beyond, Springer, 2001.

\bibitem{maclane}
S.\ Mac~Lane,
\emph{Categories for the Working Mathematician}, 2nd edition.
Springer, 1998.

\bibitem{vol5}
Vol5 source: \texttt{hott-riemann-hypothesis-vol5/lean/Vol5/YonedaNymanTrackA/}.

\bibitem{nextsteps}
Vol5 \texttt{sources/next-steps.md},
\S2.2 \emph{Rational-Dilation Externalization} and Recommended Proof
Sprint Order item~4.

\bibitem{vol5kb}
Vol5 \texttt{.knowledge-base.md} (Vol6 mirror).

\bibitem{lean4}
The Lean~4 Theorem Prover team,
\emph{Lean~4 Reference Manual}, 2024.

\bibitem{mathlib}
The mathlib Community,
\emph{The Lean Mathematical Library}.
\url{https://leanprover-community.github.io/}

\end{thebibliography}

\end{document}
